%\documentclass[12pt,oneside]{amsart}
\documentclass{scrartcl}

%\documentclass{article}

\usepackage{amsthm,amsmath,amssymb}
\usepackage[numbers]{natbib}

%\usepackage[colorlinks]{hyperref}
\usepackage{hyperref}
\hypersetup{
    colorlinks,%
    citecolor=black,%
    filecolor=black,%
    linkcolor=black,%
    urlcolor=black
}
\usepackage{hypernat}
\usepackage[top=2.5cm, bottom=2.5cm, left=2.8cm,
right=2.8cm]{geometry}
\usepackage{graphicx}


%\setlength{\parskip}{0pt}
%\setlength{\parsep}{0pt}
%\setlength{\headsep}{0pt}
%\setlength{\topskip}{0pt}
%\setlength{\topmargin}{0pt}
%\setlength{\topsep}{0pt}
%\setlength{\partopsep}{0pt}

\linespread{1}

%\usepackage{marvosym}

%\textwidth = 6.5 in \textheight = 9.2 in \oddsidemargin = 0.0 in
%\evensidemargin = 0.0 in \topmargin = -0.0 in \headheight = 0.0 in
%\headsep = 0.0in \parskip = 0.0in \parindent = 0.0in
%\def\baselinestretch{0.871}


\newtheorem{theorem}{Theorem}[section]
\newtheorem{proposition}[theorem]{Proposition}
\newtheorem{problem}[theorem]{Problem}
\newtheorem{fact}[theorem]{Fact}
\newtheorem{lemma}[theorem]{Lemma}
\newtheorem{defn}[theorem]{Definition}
\newtheorem{corollary}[theorem]{Corollary}
\newtheorem{remark}{Remark}[section]
\newtheorem{example}[theorem]{Example}

\newcommand{\E}{{\mathbb E}}
\newcommand{\se}{{\mathcal{E}}}
\newcommand{\R}{{\mathbb R}}
\renewcommand{\P}{{\mathbb P}}
\newcommand{\ac}{{\mathcal{A}}}
\newcommand{\A}{{\mathcal{A}}}
\newcommand{\B}{{\mathcal{B}}}
\newcommand{\C}{{\mathcal{C}}}
\newcommand{\M}{{\mathcal{M}}}
\newcommand{\Mf}{{\mathfrak{M}}}
\newcommand{\T}{{\mathcal{T}}}
\newcommand{\Z}{{\mathcal{Z}}}
\newcommand{\K}{{\mathcal{K}}}
\newcommand{\lc}{{\mathcal{L}}}
\newcommand{\Ps}{{\mathcal{P}}}
\newcommand{\G}{{\mathcal{G}}}
\newcommand{\pa}{{\mathring{p}}}
\newcommand{\F}{{\cal F}}
\newcommand{\samp}{{\mathcal{X}}}
\newcommand{\X}{{\mathcal{X}}}
\newcommand{\hi}{{\hat{\Phi}}}
\newcommand{\data}{{\text{data}}}
\newcommand{\relu}{{\text{ReLU}}}

\newcommand{\ridge}{{\hat{\beta}_{\text{ridge}}(\lambda)}}
\newcommand{\lasso}{{\hat{\beta}_{\text{lasso}}(\lambda)}}
\newcommand{\ridgemu}{{\hat{\mu}_t^{\text{ridge}}(\lambda)}}
\newcommand{\lassomu}{{\hat{\mu}_t^{\text{lasso}}(\lambda)}}


\newcounter{rcnt}[section]
\renewcommand{\thercnt}{(\roman{rcnt})}

\def\qt#1{\qquad\text{#1}}

\def\argmin{\mathop{\rm argmin}}
\def\argmax{\mathop{\rm argmax}}


\setlength{\parskip}{1.4 \medskipamount}
%\sloppy
%\linespread{1.3}

\begin{document}

\title{STAT 238 - Bayesian Statistics  \\
  Lecture Nineteen} 
\subtitle{Spring 2026, UC Berkeley}

\author{Aditya Guntuboyina}


\date{09 March 2026}

\maketitle

\section{A High-Dimensional Linear Regression Model}
We have a response variable $y$ and a single covariate $x$. In our
example, $y$ denotes weekly earnings and $x$ denotes years of
experience. The covariate $x$ takes the values $0, 1, \dots, m$ for
some fixed integer $m$. 

Our data is $(x_i, y_i), i = 1, \dots, n$. The model is:
\begin{align}\label{ibm_mod}
  y_i = \beta_0 + \beta_1 x_i + \beta_2 \relu(x_i - 1) + \dots +
  \beta_{m}\relu(x_i - (m-1)) + \epsilon_i
\end{align}
where $\relu(u) := \max(u, 0)$, and $\epsilon_i
\overset{\text{i.i.d}}{\sim} N(0, \sigma^2)$. Note we did not include
$\relu(x_i - m)$ because it always equals 0. 

This linear regression equation can be written in the usual notation as:
\begin{align}\label{linreg}
  y = X \beta + \epsilon
\end{align}
where $y = (y_1, \dots, y_n)^T$ and $X$ is the $n \times (m+1)$ matrix
with columns $1$, $x$, $\relu(x - j)$ for $j = 1, \dots, m-1$.

\section{Recap: Linear Regression with default priors}

As we saw before, the default prior on $\beta_0, \dots, \beta_m$ in
the linear regression model \eqref{linreg} is:
\begin{align}\label{prior_uniform}
  \beta_0, \dots, \beta_{m} \overset{\text{i.i.d}}{\sim}
  \text{uniform}(-C, C) 
\end{align}
with $C \rightarrow \infty$. With this prior, the posterior of $\beta$
given $\sigma$ becomes (see e.g., Problem 7(a) in Homework Three)
\begin{align}\label{def_post}
  \beta \mid \text{data}, \sigma \sim N\left((X^T X)^{-1} X^T y,
  \sigma^2 (X^T X)^{-1} \right)
\end{align}
which means that the posterior mean of $\beta$ is just the least
squares estimator $(X^T X)^{-1} X^T y$ (note that this posterior mean
of $\beta$ does not depend on $\sigma$, so is both the conditional
posterior mean given $\sigma$, as well as the unconditional posterior
mean).

The fact that the posterior mean of $\beta$ coincides with the least
squares estimator is shared by some priors other than
\eqref{prior_uniform} as well. For example, this is also true for the
following prior:
\begin{align}\label{def_norm}
  \beta_0, \dots, \beta_m \overset{\text{i.i.d}}{\sim} N(0, C) 
\end{align}
as $C \rightarrow \infty$. This can be seen by the following general
result (which follows, for example, from Problem 9 in Homework Two):
\begin{align}\label{normnorm}
  \beta \mid \sigma \sim N(0, Q) \implies 
  \beta \mid \text{data}, \sigma \sim N \left(\left(\frac{X^T
  X}{\sigma^2} + Q^{-1} \right)^{-1} \frac{X^T y}{\sigma^2},
  \left(\frac{X^T X}{\sigma^2} + Q^{-1}\right)^{-1} \right). 
\end{align}
From here, it is clear that if $Q$ is chosen so that $Q^{-1} \approx
0$, then we get the same posterior as \eqref{def_post}. The prior
\eqref{def_norm} corresponds to $Q = C I$ so that $Q^{-1} = (1/C)I$
which goes to zero when $C \rightarrow \infty$. So the prior
\eqref{def_norm} leads to the same posterior \eqref{def_post}. 

\section{Regularization}
In the earnings-experience regression example, we have seen in the
last class that the least squares estimator for \eqref{linreg} is not
interpretable. We want the fitted curve
\begin{align*}
  x \mapsto \hat{\beta}_0 + \hat{\beta}_1 x + \sum_{j=2}^m
  \hat{\beta}_j \relu(x - (j-1))
\end{align*}
to be smooth for it to be interpretable. This can be achieved in the
Bayesian setting if we use the prior: 
\begin{align}\label{regu_prior}
  \beta_0, \beta_1 \overset{\text{i.i.d}}{\sim} N(0, C) ~~ \text{ and
  } ~~ \beta_2, \dots, \beta_{m} \overset{\text{i.i.d}}{\sim} N(0,
  \tau^2). 
\end{align}
for some small $\tau$. The above prior
is equivalent to $\beta \sim N(0, Q)$ where $Q$ is the $(m+1) \times
(m+1)$ diagonal matrix with diagonal entries $C, C, \tau^2, \dots,
\tau^2$. So using the formula \eqref{normnorm}, the posterior mean
corresponding to this prior is:
\begin{align*}
  \E \left(\beta \mid \text{data}, \sigma, \tau \right) =
  \left(\frac{X^T X}{\sigma^2} + Q^{-1} \right)^{-1} \frac{X^T
  y}{\sigma^2}. 
\end{align*}
Because $Q$ is diagonal with diagonal entries $C, C, 1/\tau^2, \dots,
1/\tau^2$, it is easy to see that $Q^{-1}$ is also diagonal with
diagonal entries $1/C, 1/C, \tau^{-2}, \dots, \tau^{-2}$. Because $C$
is large, we can approximate $Q^{-1}$ by the matrix with diagonal
entries $0, 0, \tau^{-2}, \dots, \tau^{-2}$. In other words:
\begin{align*}
   Q^{-1} \approx \frac{1}{\tau^2} J ~~~ \text{ where } J = \begin{pmatrix}
      0 & 0 & 0 & 0 & \cdot & \cdot & \cdot & 0 \\
      0 & 0 & 0 & 0 & \cdot & \cdot & \cdot & 0 \\
      0 & 0 & 1 & 0 &\cdot & \cdot & \cdot & 0 \\
      0 & 0 & 0 & 1 & \cdot & \cdot & \cdot & 0 \\
      \cdot & \cdot & \cdot & \cdot & \cdot & \cdot & \cdot & \cdot \\
      \cdot & \cdot & \cdot & \cdot & \cdot & \cdot & \cdot & \cdot \\
      \cdot & \cdot & \cdot & \cdot & \cdot & \cdot & \cdot & \cdot \\
      0 & 0 & 0 & 0 & \cdot & \cdot & \cdot & 1 \\
\end{pmatrix}  
\end{align*}
Note that $J$ is an $(m+1) \times (m+1)$ diagonal matrix. Thus the
posterior mean becomes: 
\begin{align*}
  \E \left(\beta \mid \text{data}, \sigma, \tau \right) =
  \left(\frac{X^T X}{\sigma^2} + \frac{J}{\tau^2} \right)^{-1}
  \frac{X^T 
  y}{\sigma^2} = \left(X^T X + \frac{J}{\tau^2/\sigma^2} \right)^{-1}
  X^T y. 
\end{align*}
We use the notation $\tau^2 = \sigma^2 \gamma^2$ or equivalently
$\gamma^2 = \tau^2/\sigma^2$. This gives
\begin{align}\label{pm_norm}
  \E \left(\beta \mid \text{data}, \sigma, \tau \right) = \left(X^T X
  + \frac{J}{\gamma^2} \right)^{-1} 
  X^T y.  
\end{align}
It can be verified that this quantity $(X^TX + J/\gamma^2)^{-1} X^T y$
coincides with a ridge regularized least squares estimator. This is
shown in the next section.

\section{Connection to Ridge Regularization}
For the model \eqref{ibm_mod}, the ridge regression estimator $\ridge$
for $\beta$ is given by the minimizer of: 
\begin{equation}\label{ridge_est}
  \begin{split}
 & \sum_{i=1}^n \left(y_i - \beta_0 - \beta_1 x_i - \beta_2
    \relu(x_i-1) - \dots - \beta_{m} \relu(x_i-(m-1)) \right)^2 \\ &+
  \lambda \left(\beta_2^2 + \beta_3^2 + \dots + \beta_{m}^2
    \right). 
  \end{split}  
\end{equation}
The objective function above can also be written as 
\begin{equation*}
  \|y - X \beta\|^2 + \lambda \sum_{j=2}^{m} \beta_j^2. 
\end{equation*}
It turns out that $\ridge$ can be written in closed form using matrix
notation. To see this, note first that the gradient of the above
objective function with respect to $\beta$ is given by 
\begin{equation*}
\nabla \left(  \|y - X \beta\|^2 + \lambda \sum_{j=2}^{m} \beta_j^2
\right) =  -2 X^T y + 2 X^T X \beta + 2 \lambda 
  \begin{pmatrix}
    0 \\ 0 \\ \beta_2 \\ \cdot \\ \cdot \\ \cdot \\ \beta_{m}
  \end{pmatrix}
\end{equation*}
Recalling that $J$ is the $(m+1) \times (m+1)$  diagonal matrix with
diagonal entries $0, 0, 1, \dots, 1$, we can write 
\begin{equation*}
\nabla \left(  \|y - X \beta\|^2 + \lambda \sum_{t=2}^{n-1} \beta_j^2
\right) =  -2 X^T y + 2 X^T X \beta + 2 \lambda 
  \begin{pmatrix}
    0 \\ 0 \\ \beta_2 \\ \cdot \\ \cdot \\ \cdot \\ \beta_{n-1}
  \end{pmatrix} = -2 X^T y + 2 X^T X \beta + 2 \lambda J \beta. 
\end{equation*}
Setting this gradient equal to zero, we get 
\begin{equation*}
  -2 X^T y + 2 X^T X \beta + 2 \lambda J \beta = 0 \implies \left(X^T
    X + \lambda J \right) \beta = X^T y. 
\end{equation*}
which gives
\begin{equation}\label{ridgeformula}
  \ridge = (X^T X + \lambda J)^{-1} X^T y. 
\end{equation}
Note that $\lambda = 0$ corresponds to least squares, and $\lambda
\rightarrow \infty$ leads to linear regression. When $\lambda$ is set
to be neither very close to 0 nor very large, we should get a smooth
estimate that still gives a good fit to the data. 

Comparing \eqref{ridgeformula} to \eqref{pm_norm}, it is clear that
they are identical when $\lambda = 1/\gamma^2$. In other words, when
we use the prior \eqref{regu_prior}, then the posterior mean (given
$\tau$ and $\sigma$) coincides with ridge regularized least squares
with regularization parameter $\lambda = 1/\gamma^2 =
\sigma^2/\tau^2$.

\section{Bayesian inference for $\gamma$ and $\sigma$} 

The big advantage of Bayesian inference is that it also allows
inference on $\gamma$ and $\sigma$. We shall use the prior: 
\begin{align*}
  \log \gamma, \log \sigma \overset{\text{i.i.d}}{\sim}
  \text{uniform}(-\infty, \infty).
\end{align*}
Recall again that $\tau = \gamma \sigma$. The joint prior on all the
parameters $(\beta, \gamma, \sigma)$ is:
\begin{align*}
  f_{\beta, \gamma, \sigma}(\beta, \gamma, \sigma) \propto
  \frac{1}{\gamma \sigma} \frac{1}{\sqrt{\det Q}} \exp
  \left(-\frac{1}{2} \beta^T Q^{-1} \beta \right). 
\end{align*}
Here $Q$ is the $(m+1) \times (m+1)$ diagonal matrix with diagonal
entries $C, C, \gamma^2 \sigma^2, \dots, \gamma^2 \sigma^2$.

With this prior, we shall argue in the next lecture that the posterior
is given by:
\begin{align}\label{postQ}
  \beta \mid \text{data}, \sigma, \gamma \sim N \left(\left(X^T X +
  \gamma^{-2} J
  \right)^{-1} X^T y, \sigma^2 \left(X^T X + \gamma^{-2} J \right)^{-1} \right)
\end{align}
and
\begin{align}\label{sig_post}
  \frac{1}{\sigma^2} \mid \text{data}, \gamma \sim
  \text{Gamma}\left(\frac{n}{2} - 1, \frac{y^T y - y^T X \left(X^T X + 
 \gamma^{-2} J \right)^{-1} X^T y}{2}. 
  \right) 
\end{align}
and
\begin{align*}
  f_{\gamma \mid \text{data}}(\gamma) &\propto   \frac{\gamma^{-m}
                                        \left|X^T X + \gamma^{-2} J  
    \right|^{-1/2}}{\left(y^T y - y^T X \left(X^T X + 
 \gamma^{-2} J \right)^{-1} X^T y \right)^{(n/2) - 1}}. 
\end{align*}
With this, inference can be carried out by first taking a grid of
$\gamma$ values and computing the above posterior (on the logarithmic
scale) at the grid points. This posterior can be used to obtain
posterior samples of $\gamma$. For each sample of $\gamma$, we can
then sample $\sigma$ using the distribution \eqref{sig_post}. Given
samples from both $\gamma$ and $\sigma$, we can then sample $\beta$
using \eqref{postQ}.






%\begin{thebibliography}{99}

%\bibitem[MacKay and Peto(1995)]{MacKay1995HierarchicalDirichlet}
%D.~J.~C. MacKay and L.~C. Peto,
%\newblock ``A hierarchical Dirichlet language model,''
%\newblock \emph{Natural Language Engineering}, vol.~1,
%no.~3, pp.~289--308, 1995.

%\bibitem[Rubin(1981)]{Rubin1981BayesianBootstrap}
%D.~B. Rubin,
%\newblock ``The Bayesian bootstrap,''
%\newblock \emph{The Annals of Statistics}, vol.~9,
%no.~1, pp.~130--134, 1981.
  
%\bibitem[Fisher(1934)]{Fisher1934}
%R.~A. Fisher,
%\newblock ``Probability, likelihood and quantity of information in the logic of
%uncertain inference,''
%\newblock \emph{Proceedings of the Royal Society of London. Series~A,
%Containing Papers of a Mathematical and Physical Character}, vol.~146,
%no.~856, pp.~1--8, 1934.

%\bibitem[Efron(2010)]{Efron}
%Efron, B. (2010)
%\textit{Large-scale inference: empirical Bayes methods for estimation,
%testing, and prediction}. Cambridge University Press, 2010.

%\bibitem[Shumway and Stoffer(2017)]{ShumwayStoffer}
%Shumway, R. H. and Stoffer, D. S. (2017)
%\textit{Time Series Analysis and Its Applications: With R
%  Examples}. Springer, 4th edition. 
  
%     \bibitem[Jaynes(2003)]{Jaynes} Jaynes, E. T, (2003)
%  \textit{Probability theory: the logic of science}. Cambridge
%  University Press, 2003.
 
%  \bibitem[Sinai(1992)]{Sinai} Sinai, Y. G, (1992)
%  \textit{Probability theory: an introductory course}. Springer
%  Verlag, 1992.  

%\bibitem[{deGroot(1986)}]{DeGrootBlackwell}DeGroot, M. H. (1986)
%       A conversation with David Blackwell.
%\newblock \emph{Statistical Science}, \textbf{1}, 40--53.

%         \bibitem[Mosteller(1987)]{Mosteller} Mosteller, F, (1987)
%  \textit{Fifty challenging problems in probability with
%    solutions}. Courier Corporation, 1987.

%    \bibitem[MacKay(2003)]{MacKay} MacKay, D. J. C., (2003)
%  \textit{Information theory, inference, and learning
%  algorithms}. Cambridge University Press, 2003.

%\bibitem[{Lindley(2013)}]{Lindley}Lindley, D. V. (2013)
%   \textit{Understanding Uncertainty}. John Wiley and sons, 2013.


%\end{thebibliography}







\end{document}

%%% Local Variables:
%%% mode: latex
%%% TeX-master: t
%%% End:
